\title{Lista de Exercícios: Criptografia Assimétrica}
\author{Prof. Gabriel Rodrigues Caldas de Aquino}
\date{Compilado em: \\ \today}

\begin{document}

\maketitle

\section{Criptografia Simétrica vs. Assimétrica}
Compare criptografia simétrica e assimétrica em termos de:

\begin{enumerate}
    \item Distribuição de chaves.
    \item Aplicações de criptografia simétrica e assimétrica.
    \item Assim como a criptografia simétrica, é correto afirmar que a criptografia assimétrica pode estar suscetível a ataques de criptoanálise e força bruta?
    \item Explique por que sistemas práticos frequentemente usam ambos os tipos de criptografia.
\end{enumerate}



\section{RSA}
Para aplicações do RSA precisamos escolher dois números primos. Dado os seguintes números:
\[
    p = 17, \quad q = 11
\]

\textbf{faça}:
\begin{enumerate}
    \item Calcule \( n \) e \( \varphi(n) \).
    \item Escolha \( e = 7 \) e calcule \( d \) tal que \( e \cdot d \equiv 1 \pmod{\varphi(n)} \). Mostre os passos que você utilizou para encontrar (pode usar o algoritmo de euclides extendido ou outra forma).
    \item Encripte \( M = 88 \) usando a chave pública \((e, n)\) e mostre a decriptação com \( d \).
\end{enumerate}

\section{Segurança do RSA}

\begin{enumerate}
    \item Como a dificuldade de fatoração se relaciona com a segurança do RSA?
    \item Qual é o efeito de se utilizar um valor $n$ pequeno no RSA?
    \item Uma vez que o valor de $\varphi(n)$ seja conhecido, podemos dizer que a segurança do RSA foi comprometida? Justifique.
\end{enumerate}

\section{Diffie–Hellman}

Considere o protocolo Diffie–Hellman com os parâmetros públicos:
\[
    q = 467,\qquad \alpha = 2
\]
(ou seja, \(q\) é primo e \(\alpha\) é uma base primitiva modulo \(q\)).

\begin{enumerate}
    \item Escolham dois expoentes privados inteiros \(X_A\) e \(X_B\). Escreva e justifique os valores escolhidos.
    \item Para cada participante da comunicação, calcule os valores $Y$ públicos de cada participante:
          \[
              Y_A \equiv \alpha^{X_A} \bmod q,\qquad Y_B \equiv \alpha^{X_B} \bmod q.
          \]
          Apresente os cálculos
    \item Supondo que \(Y_A\) e \(Y_B\) sejam trocados publicamente, calcule a chave secreta comum \(K\) usando cada um dos lados e verifique que ambos obtêm o mesmo valor:
          \[
              K_A \equiv Y_B^{X_A} \bmod q,\qquad K_B \equiv Y_A^{X_B} \bmod q.
          \]
          Mostre os cálculos e verifique \(K_A = K_B\).


\end{enumerate}

\section{Segurança do Diffie-Hellman}
\begin{enumerate}
    \item Explique por que um atacante que conheça apenas \(q,\alpha,Y_A,Y_B\) não consegue, de forma prática, obter \(K\)
    \item Em aplicações práticas, qual o problema de se escolher um valor pequeno em  $q$. Justifique com base no conceito de função de mão única.
    \item É correto afirmar que o diffie-Hellman é sucetível a ataques do tipo man-in-the-middle? Descreva uma forma de mitigar esse problema caso ele exista.
\end{enumerate}

\section{Aplicações da Criptografia Assimétrica}

A criptografia assimétrica é amplamente utilizada em diversas aplicações reais de segurança da informação.


\begin{enumerate}
    \item Cite duas aplicações práticas do uso de criptografia assimétrica e justifique.
    \item Descreva o cenário em que cada uma dessas aplicações é preferível.
\end{enumerate}

Considerando aplicações de criptografia assimetrica, faça:

\begin{enumerate}

    \item  Explique o papel das chaves pública e privada nas seguintes aplicações, indicando \textbf{quem usa qual chave} e \textbf{para quê}:
          \begin{enumerate}
              \item Assinatura digital.
              \item Autenticação de usuários ou servidores (por exemplo, em SSH ou HTTPS).
              \item Confidencialidade de mensagens.
          \end{enumerate}

    \item  Descreva como os \textbf{certificados digitais} garantem a autenticidade das chaves públicas em sistemas como HTTPS. Indique:
          \begin{itemize}
              \item quem emite o certificado,
              \item o que é assinado dentro dele,
              \item e por que o navegador confia nessa assinatura.
          \end{itemize}

    \item  A criptografia assimétrica é usada na \textbf{troca de chaves} (por exemplo, no início de uma conexão TLS).
          Explique como a combinação dos dois tipos de criptografia (assimétrica + simétrica) pode ser usada em um cenário real.

    \item  Em uma frase, resuma a principal diferença entre o uso da criptografia assimétrica para \textbf{autenticação} e para \textbf{confidencialidade}.

\end{enumerate}


\end{document}